\documentclass[12pt]{article}
\usepackage{graphicx,amssymb,amsmath}
\usepackage{fullpage,times}
\usepackage{listings}
\usepackage{amsmath,mathrsfs}
\usepackage{amsmath}
\usepackage{amssymb}
\usepackage{xspace}
\usepackage{url}
\usepackage{color}
\usepackage{hyperref}

\usepackage{times}
\usepackage{amsmath,amsthm,amssymb}
\usepackage{graphics}
\usepackage{graphicx}
\usepackage{xspace}
% \usepackage[margin=32mm,nohead,nofoot]{geometry}
\usepackage[margin=20mm,nohead]{geometry}

\usepackage{framed}
\usepackage[usenames,dvipsnames]{xcolor}
\usepackage{url}
\usepackage{soul}
\usepackage{dsfont}
\usepackage{amsmath}
\usepackage{mathrsfs}
\usepackage{hyperref}
\usepackage{setspace}
\usepackage{enumitem}
\usepackage[authoryear]{natbib}
\usepackage{response-rebuttal}
\usepackage{xspace}

\newenvironment{pf}[1][Proof]{\noindent\textbf{#1.} }{\hfill\rule{1mm}{2mm}}

\newcommand*{\tool}{{\textsc{Spectre}}\xspace}

\newcommand{\code}[1]{{\fontfamily{cmtt}\fontseries{m}\fontshape{n}\selectfont\small{#1}}}

\newenvironment{shadequote}
{\begin{snugshade}\begin{quote}}
{\hfill\end{quote}\end{snugshade}}
\definecolor{shadecolor}{rgb}{0.9,0.9,0.9}

\lstset{breaklines=true,basicstyle=\itshape\footnotesize}

\begin{document}


\begin{center}
\text{\Large{Revision Response Letter}}
\end{center}


\bigskip

We would like to thank the Editor and Reviewer 3 for their valuable comments and constructive suggestions, which have been carefully taken into account when producing this revised version. 
Note that our changes and new materials are marked in \textcolor{blue}{blue} in this revised version. 
% Minor editorial changes are not highlighted for conciseness.
In this response letter, we will address the required changes summarized by the Editor and then respond to each comment raised by Reviewer 3 in detail.


\bigskip

\noindent
\textbf{\Large Required changes summarized by Editor}

\begin{shaded}
    \noindent \textbf{Comments to the Author}\\
The paper is generally recognized as making an innovative contribution within its field. Some minor issues related to writing style remain and should be addressed in a final minor revision round. Please see the comments of the third reviewer and take this opportunity to revise the overall writing style for improved clarity and fluency.   
\end{shaded}
\noindent \textbf{Response:}
Thank you for your positive feedback and helpful suggestions.
\begin{itemize}
\item We have clarified in Section I, Paragraph 3 (Page 1, right, line -2 to Page 2, left, line 7) that the theoretical proof of our work differs from that of Reference 10, and the distinction is now explicitly discussed.
\item We have carefully revised the manuscript to improve the overall writing style and have corrected all grammatical and expression issues for better clarity and fluency.
\end{itemize}
\hfill  $\blacksquare$


\bigskip
\pagebreak


\bigskip

\pagebreak

\noindent
\textbf{\Large Comments from Reviewer 3}

\begin{shaded}
    \noindent \textbf{Comment $3.1$}\\
I note that Reference 10 studied the $h$-extra $r$-component connectivity of 3-ary $n$-cube. Is the theoretical proof of this work similar to that of Reference 10?
\end{shaded}
\noindent \textbf{Response:} 
Thank you for the question. Reference 10 investigates the $h$-extra $r$-component connectivity of 3-ary $n$-cubes for $h=0,1$ and $r=3$. In contrast, our work focuses on augmented $k$-ary $n$-cubes ($AQ_{n,k}$) for general $r$, which feature a more complex dual-edge structure and richer connectivity than 3-ary $n$-cubes.
While we have carefully examined the proof techniques in Reference 10, their methods are tailored to the specific structural properties of 3-ary $n$-cubes and cannot be directly extended to augmented $k$-ary $n$-cubes. In our paper, we present new lemmas and arguments to handle the $h$-extra $r$-component connectivity of $AQ_{n,k}$. Therefore, although the overall proof framework shares some conceptual similarity, the detailed reasoning and techniques in our work are substantially different.

We have clarified this distinction in the revised manuscript (Pages 1-2 in Section I), highlighting that prior proofs are topology-specific and cannot be directly generalized to augmented $k$-ary $n$-cubes.
\hfill  $\blacksquare$

\begin{shaded}
    \noindent \textbf{Comment $3.2$}\\
2. The Grammatical mistakes need to be further improved.\\
1) Page 1, right, lines 45 and 47: ``to'' should be deleted.\\
2) Page 2, left, line 12: ``more'' should be deleted.\\
3) Page 2, right, line 31:  ``the robustness of a networks fault tolerance'' is confusing and unclear.\\
4) Page 3, right, line 11: ``adjacent to some nodes in S'' $\rightarrow$ ``adjacent to nodes of S''.\\
5) Page 3, right, line 24: ``no edges'' should be ``no edge''.\\
6) Page 3, right, line 35: ``F'' is inconsistent.\\
7) Page 5, left, line 5: ``$AQ_{n, k} - F$ is an $r$-component cut'' is not correct.\\
8) Page 5, left, line 6: ``it'' $\rightarrow$ ``$AQ_{n, k} - F$''.\\
9) Page 5, left, line 14: the expression ``as the network scale'' is not accurate.\\
10) Page 5, left, line 15: ``more'' should be deleted.\\
11) Page 6, left, lines 44 and 49: ``no isolated nodes''  $\rightarrow$ ``no isolated node''.\\
12) Page 6, left, line 49: $AQ_{n, k} - F - \{x_1, y_1, x_2, y_2\}$" $\rightarrow$ ``$AQ_{n, k} - F - \{x_1, y_1, x_2, y_2\}$''.\\
13) Page 6, left, line 57: the expression ``$i \in I_k = \{0, k - 1\}$'' is not accurate. Maybe it should be corrected to ``$i \in I_k$ and $I_k = \{0, k - 1\}$''\\
14) Page 6, right, line 3: ``we can proceed to prove in two distinct cases'' should be ``we proceed to prove two cases as below''.\\
15) Page 6, right, line 15: ``$i = \{0, 1, 2, 3\}$'' is not accurate.\\
16) Page 7, left, line 51: ``additional'' should be deleted.\\
17) Page 7, left, lines 52 and 55: ``$N(x_m, y_m)$'' should be ``$N(\{x_m, y_m\}$''.\\
18) Page 8, left, line 27: ``the algorithm 1'' should be ``Algorithm 1''.\\
19) Page 8, left, line 46: ``$o(|E|.\alpha(|V|)$'' should be ``$o(|E|.\alpha|V|))$''.\\
I hope the authors carefully check other issues of this manuscript. Thanks.

\end{shaded}
\noindent \textbf{Response:} 
We sincerely thank the reviewer for the careful proofreading and valuable suggestions. All the grammatical and expression issues listed have been carefully checked and corrected in the revised manuscript.
\hfill  $\blacksquare$



% \bibliographystyle{unsrt}
% \bibliography{TC-minor-revision-ref}

\end{document}
